\section*{Bab \ref{ch:g12:contdist} --- Peubah Acak Sinambung}

\begin{solution}{exo:g12:contdist:1}
$\P(X \leq 5) = \frac{5}{15} = \frac13$;
$\P(4 \leq X \leq 10) = \frac{6}{15} = \frac25$;
$\E(X) = \frac{15}{2} = 7.5$ menit;
$\sigma(X) = \frac{15}{\sqrt{12}} = \frac{5\sqrt3}{2} \approx 4.3$ menit.
\end{solution}

\begin{solution}{exo:g12:contdist:2}
Jelas $f \geq 0$ dan
$\int_0^2 \frac38 t^2\,\dd t = \frac38\left[\frac{t^3}{3}\right]_0^2
= \frac38 \times \frac83 = 1$: jadi sebuah \omterm{def:g12:contdist:density}{kerapatan}.
\[
\E(X) = \int_0^2 \frac38 t^3\,\dd t = \frac38 \times \frac{16}{4} = \frac32,
\qquad
\P(X \geq 1) = \int_1^2 \frac38 t^2\,\dd t = \frac{8 - 1}{8} = \frac78 .
\]
\end{solution}

\begin{solution}{exo:g12:contdist:3}
\emph{1.} $\E(X) = \frac{1}{0.2} = 5$ tahun;
$\P(X > 5) = \eu^{-0.2\times5} = \eu^{-1} \approx 0.37$.

\emph{2.} Menurut ketiadaan ingatannya (\cref{thm:g12:contdist:memoryless}),
$\pcond{X > 3}{X > 8} = \P(X > 5) = \eu^{-1} \approx 0.37$: jadi ketiga tahun
pemakaiannya tidak mengubah apa pun.
\end{solution}

\begin{solution}{exo:g12:contdist:4}
$\P(X > m) = \eu^{-\lambda m} = \frac12$ memberi
$m = \frac{\ln 2}{\lambda} \approx \frac{0.69}{\lambda}$, yang lebih kecil daripada
$\E(X) = \frac1\lambda$. \omterm{def:g12:contdist:density}{Kerapatan} eksponennya berekor kanan yang panjang:
beberapa masa hidup yang luar biasa panjang menarik \emph{rata-ratanya} ke atas
\emph{\omterm{def:g11:stat:median}{mediannya}}, yaitu nilai yang dilampaui separuh populasinya. (Inilah
waktu paruh pada \cref{ex:g12:exp:halflife}: separuh atomnya melewatinya, walaupun
rata-rata masa hidupnya lebih panjang.)
\end{solution}

\begin{solution}{exo:g12:contdist:5}
Di sini $168 = \mu - \sigma$ dan $182 = \mu + \sigma$: jadi kira-kira $68\%$.
Lalu $189 = \mu + 2\sigma$: di atasnya terletak separuh dari sisanya,
$100 - 95.4 = 4.6\%$, jadi kira-kira $2.3\%$.
Lalu $154 = \mu - 3\sigma$: kira-kira $\frac{100 - 99.7}{2} = 0.15\%$.
\end{solution}

\begin{solution}{exo:g12:contdist:6}
$\P(X < 500) \leq 2.5\%$ berarti $500$ harus duduk sekurang-kurangnya dua \omterm{def:g11:stat:variance}{simpangan
baku} di bawah rata-ratanya (\omterm{def:g12:contdist:normal}{distribusi normalnya} meninggalkan $2.3\% \approx 2.5\%$ di bawah
$\mu - 2\sigma$; dengan $1.96$ yang lazim, penalarannya identik):
\[
\mu - 2\sigma \geq 500 \iff \mu \geq 500 + 8 = 508 \text{ g}.
\]
Jadi mesinnya harus disetel ke $508$ g rata-rata --- dan harga
jaminannya adalah $8$ g produk ``gratis'' tiap kantongnya.
\end{solution}

\begin{solution}{exo:g12:contdist:7}
\emph{1.} Dengan $p = \frac16$, $n = 1200$:
$\sqrt{\frac{p(1-p)}{n}} = \sqrt{\frac{5/36}{1200}} \approx 0.0108$, sehingga
selang ayunan $95\%$-nya adalah
\[
\frac16 \pm 1.96 \times 0.0108 \approx \intcc{0.146}{0.188}.
\]

\emph{2.} Frekuensi yang teramati $\frac{260}{1200} \approx 0.217$ terletak jauh
di luarnya: jadi hipotesis kesetimbangannya ditolak pada taraf $5\%$.
Simpangannya kira-kira $4.6$ \omterm{def:g11:stat:variance}{simpangan baku} --- yang bagi \omterm{def:g10:numbers:approx}{hampiran}
normal berpeluang berorde $10^{-6}$, jauh lebih meyakinkan daripada
batas $\leq 4.6\%$ dari Bienaymé--Chebyshev
(\cref{exo:g12:sums:7}).
\end{solution}

\begin{solution}{exo:g12:contdist:8}
\emph{1.} Marjinnya $\frac{1}{\sqrt n} \leq 0.02$ menuntut
$n \geq 2500$ orang.

\emph{2.} Untuk membedakan perolehan yang berjarak $1$ angka, marjinnya harus jauh
di bawah $0.5$ angka, sehingga menuntut $n \geq \left(\frac{1}{0.005}\right)^2 =
40\,000$ menurut rumus yang disederhanakan --- dan itu sudah tak praktis bagi kebanyakan jajak pendapat.
Lebih buruk lagi, marjin statistiknya hanya memperhitungkan galat \emph{pencuplikan};
sedangkan bias sistematisnya (contoh yang tak mewakili, penolakan menjawab, perubahan
menit terakhir) tidak menyusut ketika $n$ bertambah dan biasanya melampaui satu angka. Jadi
perlombaan berjarak $1$ angka memang terlalu ketat untuk ditetapkan.
\end{solution}

\begin{solution}{exo:g12:contdist:9}
\emph{1.} Karena $\alpha > 0$, berlaku $c \leq \alpha X + \beta \leq d
\iff \frac{c - \beta}{\alpha} \leq X \leq \frac{d - \beta}{\alpha}$, sehingga
\[
\P(c \leq Y \leq d)
= \int_{(c-\beta)/\alpha}^{(d-\beta)/\alpha} f(t)\,\dd t .
\]
Substitusi $y = \alpha t + \beta$ (yaitu membaca luasnya dalam peubah $y$,
dengan $t = \frac{y - \beta}{\alpha}$, $\dd t = \frac{\dd y}{\alpha}$)
mengubahnya menjadi
$\int_c^d \frac{1}{\alpha} f\left(\frac{y-\beta}{\alpha}\right)\dd y$: jadi
\omterm{def:g11:func:function}{fungsi} $g(y) = \frac1\alpha f\left(\frac{y-\beta}{\alpha}\right)$ adalah
\omterm{def:g12:contdist:density}{kerapatan} bagi $Y$.

\emph{2.} Dengan $f = \varphi$, $\alpha = \sigma$, $\beta = \mu$:
\[
g(y) = \frac{1}{\sigma}\,\varphi\!\left(\frac{y - \mu}{\sigma}\right)
= \frac{1}{\sigma\sqrt{2\pi}}\,
\exp\left(-\frac{(y-\mu)^2}{2\sigma^2}\right),
\]
yang karena itu merupakan \omterm{def:g12:contdist:density}{kerapatan}
$\mathcal N(\mu, \sigma^2) = \mu + \sigma\,\mathcal N(0,1)$.
\end{solution}

\begin{solution}{pb:g12:contdist:1}
\textbf{1.} $\P(\text{tunggu} \leq 10) = \frac{10}{30} =
\frac13$; $\E = 15$ menit;
$\sigma = \frac{30}{\sqrt{12}} \approx 8.7$ menit.

\textbf{2.} $\int_0^1 2x\,\dd x = 1$: jadi sebuah \omterm{def:g12:contdist:density}{kerapatan}.
$\E(X) = \int_0^1 2x^2\,\dd x = \frac23$;
$\P\left(X > \frac12\right) = \int_{1/2}^1 2x\,\dd x =
1 - \frac14 = \frac34$.

\textbf{3.} $\P(T > 15) = \eu^{-15/10} \approx 0.22$. \omterm{def:g10:stats:median}{Mediannya}:
$\eu^{-m/10} = \frac12$, jadi $m = 10\ln 2 \approx 6.9$ menit ---
lebih kecil daripada rata-ratanya $10$ sebab ekor kanan eksponennya
yang panjang (yaitu penungguan raksasa yang jarang) menyeret rata-ratanya ke atas sambil membiarkan
\omterm{def:g11:stat:median}{mediannya} pada penungguan yang ``lazim''.

\textbf{4.} Menurut ketiadaan ingatannya,
$\P(T > 35 \mid T > 20) = \P(T > 15) \approx 0.22$: jadi
kedua puluh menit yang sudah dijalani itu tak membeli apa-apa --- arusnya
tidak menua.

\textbf{5.} (a) eksponen (sebab peluruhan radioaktif adalah model yang
untuknya ketiadaan ingatan diciptakan); (b) seragam pada
$\intcc{0}{8}$ (yaitu jadwal dengan pergeseran acak); (c)
eksponen; (d) bukan keduanya: bola lampu yang aus \emph{lebih} mungkin
putus pada jam berikutnya daripada yang baru, sehingga
$\P(T > s + t \mid T > s) < \P(T > t)$: jadi penuaannya bertentangan dengan
teorema ketiadaan ingatannya.

\textbf{6.} Berturut-turut $68\,\%$, $95\,\%$, $99.7\,\%$: yaitu ketiga
selang tetengernya.

\textbf{7.} $z = \frac{190 - 172}{8} = 2.25$, jadi ekor atasnya
$\approx 1.2\,\%$.

\textbf{8.} Di sini $z = 2$: jadi kira-kira $2.3\,\%$ --- kurang lebih satu orang
dari $44$.

\textbf{9.} Syaratnya duduk pada $\pm 2\sigma$: jadi kira-kira
$95.4\,\%$ lolos dan $4.6\,\%$ ditolak --- dan itu membangkrutkan pada skala besar. Pada
$\pm 6\sigma$ laju kegagalannya jatuh menjadi kira-kira dua bagian dari tiap
\emph{miliar}: jadi enam sigma membeli hak untuk memproduksi massal
tanpa memeriksa secara massal.

\textbf{10.} Di sini $\sigma = 5$; dan dengan koreksi kesinambungannya,
$\P \approx \P\left(\abs Z \leq \frac{5.5}{5}\right) =
\P(\abs Z \leq 1.1) \approx 0.73$. Teoremanya menghaluskan
papan Galton pada \cref{pb:g11:binom:1}: tangga lubangnya
menjadi lonceng yang sinambung.

\textbf{11.} $n \approx \frac{1.96^2 \times 0.25}{0.005^2}
\approx 38\,400$ orang --- dan pada ukuran itu galat
\emph{di luar} pencuplikannya (kerangkanya, penolakannya, kebohongannya) menenggelamkan
marjinnya: jadi perlombaan satu angka berada di luar jangkauan jajak pendapat yang jujur, dan
itulah sebabnya lembaga yang serius lalu berkata ``terlalu ketat untuk ditetapkan''.

\textbf{12.} Ketiadaan ingatannya: sejak saat kamu tiba,
sisa penungguannya eksponen dengan rata-rata $10$ --- yaitu $10$
menit penuh, bukan $5$. Kata ``jarak rata-rata $10$'' pada jadwalnya
menyesatkan: penungguanmu berdistribusi sama dengan satu jarak yang utuh.

\textbf{13.} Kedatangan yang acak mendarat sebanding dengan panjang
jaraknya: $\P(\text{jarak panjang}) = \frac{15}{20} = \frac34$.
Harapan penungguannya: $\frac34 \times \frac{15}{2} + \frac14 \times
\frac52 = 5.625 + 0.625 = 6.25$ menit --- lebih daripada nilai naif
$5$, walaupun jarak rata-ratanya $10$. Biang keladinya: \emph{pencuplikan
yang berbias panjang} --- sebab jarak yang panjang menangkap lebih banyak penumpang.

\textbf{14.} Penungguan mundur dan majunya masing-masing eksponen
dengan rata-rata $10$ (sebab ketiadaan ingatannya berjalan ke dua arah dari
kedatanganmu): jadi jarak yang memuatmu berpanjang rata-rata $20$ menit
--- dua kali jarak yang lazim. Paradoks pemeriksaan dalam satu kalimat:
\emph{selang yang kebetulan kamu periksa bukanlah selang
yang lazim, sebab kamu lebih mungkin jatuh ke dalam selang yang
besar.}

\textbf{15.} Kelas rata-ratanya: $\frac{9 \times 10 + 110}{10} =
20$ murid. \omterm{def:g11:stat:mean}{Rata-rata} yang dialami muridnya:
$\frac{90 \times 10 + 110 \times 110}{200} = \frac{13\,000}
{200} = 65$ murid. Brosurnya mencetak $20$; padahal tiga dari lima
muridnya duduk di kelas berisi $110$ dan menjalani angka $65$ itu.

\textbf{16.} Orang yang populer muncul pada banyak daftar teman, sehingga
mencuplik ``seorang teman'' berbias panjang ke arah yang pandai bergaul
--- jadi temanmu diambil dari jarak yang panjang pada jadwal
pergaulannya.

\textbf{17.} Untuk $t \geq 0$:
$\P(T > t) = \P(-10\ln U > t) = \P\left(U <
\eu^{-t/10}\right) = \eu^{-t/10}$: jadi tepat \omterm{def:g11:func:function}{fungsi}
kesintasan eksponennya --- satu \omterm{def:g12:exp:ln}{logaritma} mengubah derau
seragam milik komputer itu menjadi waktu tunggu apa pun.

\textbf{18.} Tiap peubah seragamnya berata-rata $\frac12$ dan beragam
$\frac{1}{12}$: jadi jumlah dua belasnya berata-rata $6$ dan beragam
$1$ --- sehingga resepnya menghasilkan rata-rata $0$ dan \omterm{def:g12:randvar:exp}{ragam} $1$. Dua belas
dorongan kecil yang \omterm{def:g12:sums:indep}{saling bebas} dijumlahkan: itulah mekanisme Galton, dan menurut
berkah De Moivre--Laplace histogramnya sudah
berbentuk lonceng di mata.

\textbf{19.} Tentang modelnya: pasar bukanlah jumlah banyak
sebab kecil yang \emph{\omterm{def:g12:sums:indep}{saling bebas}} --- kepanikan mengorelasikan
segalanya (yaitu pelajaran tahun 2008 pada soal sebelumnya) dan
menghasilkan ekor gemuk yang tak dimiliki \omterm{def:g12:contdist:density}{kerapatan} normal mana pun. Ketika datanya
membisikkan $10^{-137}$, katekismus sang ahli statistika menjawab:
\emph{modelnyalah yang ditolak, bukan harinya}.

\textbf{20.} Seragam: semua kedudukan setara, yaitu kejujuran ``aku
hanya tahu batasnya''. Eksponen: bahaya yang tak pernah
menua --- peluruhan, kedatangan, bunyi klik. Normal: lonceng
demokratis dari banyak sebab kecil yang \omterm{def:g12:sums:indep}{saling bebas} --- tinggi badan, galat,
rata-rata. Sengatannya: apa yang kamu cuplik terbias oleh cara kamu tersandung
padanya --- bus, kelas, teman. Dan busurnya: anak
yang mencacah kardus jus di kelas 6, sesudah sepuluh jilid
soal, kini mengukur kerumunan dengan sebuah kurva; sedangkan epsilon,
\omterm{def:g12:integ:area}{integral} dan teorema limit pusat menunggu di
jilid universitas untuk menjelaskan \emph{mengapa} loncengnya berdentang bagi
segalanya.

\end{solution}
